\chapter{The volatility smile and the Heston model}
\label{cap:heston}

This chapter presents the Heston (1993) model: the phenomenon it must capture (the \emph{skew}), the closed-form
characteristic function, and two independent Fourier engines, COS and Lewis. Monte Carlo is covered in
Chapter~\ref{cap:mc}.

\section{Why Black is not enough: smile and skew}
If the Black model were correct, inverting the market prices of all options of one expiry would always return the
same volatility. On the SPX index, this is what happens instead:
\begin{itemize}
\item implied volatility is \textbf{decreasing in $K$} (\emph{skew}): OTM puts are more expensive than OTM calls;
\item the skew is \textbf{steeper at short expiries} and flattens as $T$ grows;
\item the ATM level also varies with expiry (\emph{term structure}).
\end{itemize}
The function $(K,T)\mapsto\sigma_{\text{impl}}(K,T)$ is called the \textbf{implied-volatility surface}.

\begin{perche}[title={Why is there a skew?}]
Recalling Breeden--Litzenberger \eqref{eq:breeden}, a negative skew means that the market density has a
\textbf{heavier left tail} than the lognormal: crashes are more likely than Black says.

The economic mechanism is the \emph{leverage effect}: when the index falls, volatility rises. We therefore need a
random volatility that is \emph{negatively correlated} with the price.
\end{perche}

\section{The Heston model}
Heston makes the instantaneous variance random and correlates it with the price. It is the simplest model that
produces the skew and still admits prices in almost closed form: the variance reverts to a long-run level, the
negative correlation generates the leverage effect, and the \emph{affine} structure of the equations makes the
characteristic function of the log-price known in closed form. With only five parameters the model describes the
whole surface, that is, all expiries at once.

\subsection{The dynamics}
We work under the $T$-forward measure $\Q^T$, as in Chapter~\ref{cap:black}. With deterministic rates and carry, the
risk-neutral Heston dynamics carry over unchanged to the forward $F_t=F(t,T)$:
\begin{equation}\label{eq:heston}
\begin{aligned}
\mathrm{d}F_t&=\sqrt{v_t}\,F_t\dd W^F_t,\\
\mathrm{d}v_t&=\kappa(\theta-v_t)\dd t+\xi\sqrt{v_t}\dd W^v_t,\qquad \dd\langle W^F,W^v\rangle_t=\rho\dd t .
\end{aligned}
\end{equation}
The instantaneous \emph{variance} $v_t$ follows a Cox--Ingersoll--Ross (CIR) process, that is, a mean-reverting
diffusion with volatility proportional to $\sqrt{v_t}$. The five parameters are summarized in
Table~\ref{tab:heston-parametri}.

\begin{table}[htbp]
\centering
\small
\caption{The parameters of the Heston model and their effect on the smile.}
\label{tab:heston-parametri}
\begin{tabularx}{\textwidth}{@{}>{$}c<{$}l>{\raggedright\arraybackslash}X@{}}
\toprule
\text{par.} & name & effect on the smile \\
\midrule
v_0 & initial variance & ATM level of short expiries \\
\theta & long-run variance & ATM level of long expiries \\
\kappa & mean-reversion speed & how fast $v$ moves from $v_0$ to $\theta$; how fast the skew flattens with $T$ \\
\xi & volatility of variance (\emph{vol-of-vol}) & fatter tails on both sides: curvature of the smile \\
\rho & price--variance correlation & slope: $\rho<0$ creates the negative skew \\
\bottomrule
\end{tabularx}
\end{table}

The required domain is $\kappa-\rho\xi>0$, automatic for $\rho\le0$ (the SPX case). The code rejects other inputs:
for $\kappa-\rho\xi\le0$ some expressions in the characteristic function become $0/0$.

\subsection{Expected variance and integrated variance}
Taking the expectation of the $v$ equation (the $\mathrm{d}W$ term has zero mean) gives an ordinary differential
equation, $\frac{\mathrm{d}}{\mathrm{d}t}\E v_t=\kappa(\theta-\E v_t)$, with solution
\begin{equation}\label{eq:var-attesa}
\bar v(t):=\E v_t=\theta+(v_0-\theta)e^{-\kappa t}.
\end{equation}
Integrating over $[0,T]$ and setting $\tau_\kappa=(1-e^{-\kappa T})/\kappa$:
\begin{equation}\label{eq:sigma-bar}
\bar\Sigma:=\int_0^T\bar v(t)\dd t=\theta T+(v_0-\theta)\,\tau_\kappa .
\end{equation}
$\bar\Sigma$ is the ``expected total variance''. The effective volatility $\sigma_{\text{eff}}=\sqrt{\bar\Sigma/T}$
will come back in the limit $\xi\to0$.

\subsection{The Feller condition}
The Feller ratio $\nu_F=2\kappa\theta/\xi^2$ decides whether the variance can reach zero.
\begin{itemize}
\item If $\nu_F\ge1$, the drift $\kappa\theta$ near 0 is strong enough to prevent it.
\item If $\nu_F<1$, the variance hits 0 and immediately leaves it (the boundary is instantaneously reflecting).
\end{itemize}
\begin{attenzione}
Realistic SPX calibrations typically \emph{violate} Feller: enough curvature requires a large $\xi$. This does
\textbf{not} cause problems for the Fourier methods of this chapter, because the characteristic function is the same
whether or not zero is attainable. It does, however, cause problems for Monte Carlo (Chapter~\ref{cap:mc}).

Imposing Feller in calibration would spoil the fit of the short-dated skew: in the project the ratio $\nu_F$ is
\emph{reported}, not imposed.
\end{attenzione}

\section{The Heston characteristic function}
In Heston the density of $F_T$ is not known in closed form, but the characteristic function is, because the model is
\textbf{affine}: the drift and the instantaneous variance are linear functions of the state $(x,v)$.

\subsection{Step-by-step derivation}
Let $X=\ln(F_T/F)$ be the log-return of the forward and $\varphi(u)=\E^T[e^{iuX}]$ its characteristic function (CF).
In the state $x_t=\ln(F_t/F)$, by Itô, $\mathrm{d}x=-\tfrac12v\dd t+\sqrt v\dd W^F$.

\textbf{Step 1 (Feynman--Kac).} The function $f(t,x,v)=\E[e^{iuX}\mid x_t=x,v_t=v]$, evaluated along the process, is a
martingale. As in \eqref{eq:feynman-kac}, its $\mathrm{d}t$ term must vanish:
\begin{equation}\label{eq:heston-pde}
f_t+\tfrac12v(f_{xx}-f_x)+\kappa(\theta-v)f_v+\tfrac12\xi^2vf_{vv}+\rho\xi vf_{xv}=0,\qquad f(T,x,v)=e^{iux}.
\end{equation}
The cross term $\rho\xi v f_{xv}$ comes from the covariation $\mathrm{d}\langle x,v\rangle=\rho\xi v\dd t$.

\textbf{Step 2 (affine ansatz).} Since the coefficients are linear in $v$, we look for an exponential-affine solution:
\[
f=\exp\big(iux+C(\tau)+vD(\tau)\big),\qquad \tau=T-t,\qquad C(0)=D(0)=0 .
\]
The derivatives are simple: $f_x=iuf$, $f_{xx}=-u^2f$, $f_v=Df$, $f_{vv}=D^2f$, $f_{xv}=iuDf$, $f_t=-(C'+vD')f$.

\textbf{Step 3 (separating the terms).} Substituting and dividing by $f$, the equation becomes linear in $v$. It must
hold for every $v$, so the constant term and the coefficient of $v$ vanish separately:
\begin{equation}\label{eq:riccati}
C'=\kappa\theta D,\qquad D'=\tfrac12\xi^2D^2-\beta D-\tfrac12A,
\qquad\beta:=\kappa-i\rho\xi u,\quad A:=u^2+iu .
\end{equation}
The second is a \textbf{Riccati} equation (quadratic in $D$) and is solved in closed form; the first is then
integrated directly.

\textbf{Step 4 (solving the Riccati equation).} The polynomial $\tfrac12\xi^2D^2-\beta D-\tfrac12A$ has roots
\[
r_\pm=\frac{\beta\pm d}{\xi^2},\qquad d=\sqrt{\beta^2+\xi^2A}\quad(\mathrm{Re}\,d>0).
\]
Separating variables,
\[
\frac{\mathrm{d}D}{(D-r_+)(D-r_-)}=\tfrac12\xi^2\dd\tau ,
\]
and integrating by partial fractions:
\[
\frac1d\ln\frac{D-r_+}{D-r_-}=\tau+\frac1d\ln\frac{r_+}{r_-}.
\]
Solving for $D$ and setting $g=r_-/r_+=(\beta-d)/(\beta+d)$ gives the ``little trap'' form of Albrecher, Mayer,
Schoutens and Tistaert (2007):
\begin{align}
D(T)&=\frac{\beta-d}{\xi^2}\,\frac{1-e^{-dT}}{1-g\,e^{-dT}},\label{eq:heston-D}\\
C(T)&=\frac{\kappa\theta}{\xi^2}\Big[(\beta-d)T-2\operatorname{Log}\frac{1-g\,e^{-dT}}{1-g}\Big],\label{eq:heston-C}\\
\varphi(u)&=e^{C+v_0D}.\label{eq:heston-cf}
\end{align}

\textbf{Step 5 (immediate checks).} For $u=0$ and for $u=-i$ we have $A=0$, hence $D=C=0$ and $\varphi=1$.
\begin{itemize}
\item $\varphi(0)=1$: it is a probability distribution.
\item $\varphi(-i)=\E[F_T/F]=1$: the forward is a martingale.
\end{itemize}

\begin{perche}[title={Why the complex square root does not jump}]
On the real axis and on the line $\mathrm{Im}\,u=-\tfrac12$ (used by Lewis) we have, respectively,
\[
\mathrm{Re}\,d^2=\kappa^2+(1-\rho^2)\xi^2u^2>0,\qquad
\mathrm{Re}\,d^2=(\kappa-\tfrac12\rho\xi)^2+(1-\rho^2)\xi^2u^2+\tfrac14\xi^2>0 .
\]

In both cases $d^2$ never crosses the negative real half-axis, where the principal square root has its branch cut.
Hence $d$ is continuous along the contours we use.
\end{perche}

\subsection{The ``little Heston trap''}
The original Heston (1993) formula contains $e^{+dT}$ and $g_1=1/g$: it is obtained from the previous one by the
substitution $d\to-d$. Algebraically it is identical to \eqref{eq:heston-D}--\eqref{eq:heston-C}, because the
solution of the Riccati equation is unique and even in $d$. Numerically it is not.

\begin{enumerate}
\item The computer's complex logarithm (principal $\operatorname{Log}$) always returns an argument in $(-\pi,\pi]$: it
has a \emph{branch cut} along the negative real half-axis.
\item In the original form the argument of the logarithm behaves like $-g_1e^{dT}/(1-g_1)$, whose phase grows like
$T\,\mathrm{Im}\,d$: as $u$ grows it spirals around the origin.
\item Every time it crosses the cut, $\operatorname{Log}$ jumps by $2\pi i$ and $\varphi$ is multiplied by
$e^{\mp4\pi i\kappa\theta/\xi^2}$, which differs from 1 unless $\nu_F$ is an integer.
\end{enumerate}
For long expiries the prices thus become \emph{silently} wrong. An example from the literature (Albrecher et
al.\ 2007): a 10-year ATM call on the FO set, 29.83 instead of 33.84.

In the little trap form, by contrast, $|e^{-dT}|<1$ and the argument of the logarithm tends to $1/(1-g)$ as
$T\to\infty$: it never winds around zero. The original form also has a second flaw: $e^{dT}$ overflows for large
$uT$.

\subsection{The form used in the code: no subtractions, \code{log1p}}
Equation \eqref{eq:heston-D} contains $\beta-d$, a difference of two nearly equal numbers when $\xi$ is small:
significant digits are lost. Moreover, the logarithm in \eqref{eq:heston-C} is $\operatorname{Log}(1+O(\xi^2))$ and is
multiplied by $2\kappa\theta/\xi^2$, which amplifies its rounding error.

\textbf{Step 1.} The subtraction is removed with the identity $(\beta-d)(\beta+d)=\beta^2-d^2=-\xi^2A$:
\[
\beta-d=-\frac{\xi^2A}{\beta+d}.
\]
\textbf{Step 2.} With $e_1=1-e^{-dT}$ we define
\[
q=\frac{\xi^2A\,e_1}{2d(\beta+d)},\qquad\text{so that}\qquad \frac{1-g\,e^{-dT}}{1-g}=1-q .
\]
\textbf{Step 3.} Equations \eqref{eq:heston-D}--\eqref{eq:heston-C} become
\begin{equation}\label{eq:heston-cf-code}
D=-\frac{A\,e_1}{2d(1-q)},\qquad
C=-\frac{\kappa\theta AT}{\beta+d}-\frac{2\kappa\theta}{\xi^2}\operatorname{Log}(1-q).
\end{equation}
\textbf{Step 4.} $\operatorname{Log}(1-q)$ is computed in real arithmetic. Writing $-q=p+is$:
\[
\operatorname{Log}(1+p+is)=\tfrac12\,\mathrm{log1p}(2p+p^2+s^2)+i\,\mathrm{atan2}(s,1+p).
\]
The function \code{log1p} computes $\ln(1+y)$ with full relative precision even for tiny $y$. In addition, $e_1$ is
computed with \code{expm1}.

\begin{progetto}
\code{src/qpricer/models/heston.py}, function \code{cf}, implements exactly \eqref{eq:heston-cf-code}. It is the same
function as the little trap form, with the same principal branch, but without subtractions.
\begin{itemize}
\item At $\xi=10^{-6}$ the CF agrees with Black's to within $10^{-10}$.
\item An independent check integrated the Riccati equation \eqref{eq:riccati} numerically with its own Runge--Kutta
scheme (no closed form, no branch choice). The discrepancy is $\le{1.1\times10^{-12}}$ in absolute value, up to $T=10$,
including extreme cases ($\xi=1.5$, $\rho=+0.5$, $\rho=-0.99$, $\nu_F=0.08$).
\item The original Heston form is kept in the tests as a \emph{negative control}: on the FO set at $T=10$ it must fail.
\end{itemize}
\end{progetto}

\section{Fourier pricing, method 1: COS}
\subsection{The idea}
The price of a put is an integral of the payoff against the density $f$ of $X$:
\[
P=\DF\int_{\R}(K-Fe^y)^+f(y)\dd y .
\]
We do not know $f$, but we can approximate it with a \emph{cosine series} on an interval $[a,b]$ that holds almost all
of the probability. The method is due to Fang and Oosterlee (2008), in the forward-centered variant of Le Floc'h
(2020).

\textbf{Step 1.} On $[a,b]$ every smooth function can be expanded in a cosine series:
\[
f(y)=\sum_{k\ge0}{}'A_k\cos(u_k(y-a)),\qquad u_k=\frac{k\pi}{b-a},
\qquad A_k=\frac2{b-a}\int_a^bf(y)\cos(u_k(y-a))\dd y .
\]
The prime on $\sum'$ indicates that the $k=0$ term is halved.

\textbf{Step 2.} The coefficient is read off the characteristic function. Since $\cos z=\mathrm{Re}\,e^{iz}$,
\[
A_k=\frac2{b-a}\,\mathrm{Re}\Big[\int_a^be^{iu_k(y-a)}f(y)\dd y\Big]
\approx\frac2{b-a}\,\mathrm{Re}\big[\varphi(u_k)e^{-iu_ka}\big].
\]
The approximation consists in extending the integral to the whole of $\R$: the mass outside $[a,b]$ is neglected.

\textbf{Step 3.} Sum and integral are swapped. The price becomes a sum of products (coefficient read off
$\varphi$)$\times$(integral of the payoff against a cosine). The second factor is computed by hand. For the put, with
$z=\ln(K/F)$ the point where the payoff vanishes:
\begin{gather*}
V_k=\frac2{b-a}\Big[K\psi_k(\bar z)-F\chi_k(\bar z)\Big],\qquad \bar z=\min(\max(z,a),b),\\
\psi_k(z)=\int_a^z\cos(u_k(y-a))\dd y,\qquad\chi_k(z)=\int_a^ze^y\cos(u_k(y-a))\dd y .
\end{gather*}
Both integrals are elementary:
\begin{gather*}
\chi_k(z)=\frac{e^{z}\cos(u_k(z-a))+u_ke^{z}\sin(u_k(z-a))-e^{a}}{1+u_k^2},\\
\psi_k(z)=\frac{\sin(u_k(z-a))}{u_k}\quad(k\ge1),\qquad\psi_0(z)=z-a .
\end{gather*}
\textbf{Step 4.} The price is
\begin{equation}\label{eq:cos}
P=\DF\sum_{k=0}^{N-1}{}'\,w_kV_k,\qquad w_k=\mathrm{Re}\big[\varphi(u_k)e^{-iu_ka}\big],
\end{equation}
where the prime on $\sum'$ indicates that the $k=0$ term is taken with weight $\tfrac12$ (since $\varphi(0)=1$, we
have $w_0=1$ and the weighted term is $\tfrac12V_0$).
Truncating $z$ to $[a,b]$ (that is, $\bar z$) is not a detail: without it, strikes outside the interval give wrong
prices, even negative ones.

\begin{perche}[title={Why COS is fast, and why it prices puts}]
\begin{itemize}
\item The Heston density is smooth and its CF is analytic: the coefficients $w_k$ decay \emph{exponentially} in $k$,
at the rate of $|\varphi(u_k)|$. A few hundred terms are enough.
\item The interval is for $X$ and does not depend on the strike: $\varphi(u_k)$ is computed once per expiry and serves
all strikes. The price of $M$ strikes is an $(M\times N)$ matrix--vector product.
\item The put payoff is bounded (at most $K$). The call payoff grows like $e^y$, and its coefficients $V_k$ are
proportional to $e^b$: every truncation error is amplified. One therefore computes the put and obtains the call by
parity: \mbox{$C=P+\DF\,(F-K)$}.
\end{itemize}
\end{perche}

\subsection{The truncation interval}
We choose an interval centered on the mean of $X$ and $L$ standard deviations wide:
\begin{equation}\label{eq:cos-range}
[a,b]=\big[c_1-L\sqrt{c_2},\ c_1+L\sqrt{c_2}\big],\qquad c_1=-\tfrac12\bar\Sigma,\quad c_2=\Var X .
\end{equation}
$c_1$ and $c_2$ are the first two cumulants of $X$. They are obtained by differentiating the cumulant generating
function $\mathcal K(w)=\ln\E e^{wX}$ (the CF evaluated at $u=-iw$) and solving the linear equations satisfied by the
coefficients of the expansion of $D$ in powers of $w$. With $E=e^{-\kappa T}$:
\begin{equation}\label{eq:c2}
\begin{split}
c_2&=2\theta\Big[\alpha_0(T-\tau_\kappa)+\alpha_1(\tau_\kappa-TE)+\tfrac12\alpha_2\kappa\tau_\kappa^2\Big]\\
&\quad+2v_0\Big[\alpha_0\tau_\kappa+\alpha_1TE+\alpha_2E\tau_\kappa\Big],
\end{split}
\end{equation}
with
\[
\alpha_0=\frac12-\frac{\rho\xi}{2\kappa}+\frac{\xi^2}{8\kappa^2},\qquad
\alpha_1=\frac{\rho\xi}{2\kappa}-\frac{\xi^2}{4\kappa^2},\qquad
\alpha_2=\frac{\xi^2}{8\kappa^2}.
\]
Check: as $T\to0$, $c_2\to v_0T$ (the variance of one Black step with volatility $\sqrt{v_0}$). For
$\kappa T<10^{-3}$ the first-order expansion in $\kappa$ is used, because the exact formula has a $1/\kappa^3$
prefactor on a numerator of order $\kappa^3$ and becomes noisy.

\begin{attenzione}[title={A misprint in the literature}]
The $c_2$ formula in Table 11 of Fang--Oosterlee (2008) contains a typo: the term $\bar u(6e^{-\lambda T}-7)$ should
read $\bar u(4e^{-\lambda T}-5)$. The printed version is too small by $\xi^2\theta(1-E)/(4\kappa^3)$ and can become
negative: this is why the authors use $|c_2|$.

In the project $c_2$ is derived in two independent ways. It is then checked against finite differences of
$\mathcal K(w)$ and against a numerical integration of the moment equations (agreement $\le{2.7\times10^{-15}}$
relative). Hand estimate on the FO set at $T=1$: $c_2=0.03157$, against $0.03091$ with the printed formula.
\end{attenzione}

\subsection{The two sources of COS error}
The COS error has two independent components.
\begin{itemize}
\item \textbf{Series error}: the sum is truncated at $N$ terms. It decreases exponentially with $N$.
\item \textbf{Range error} (or tail error): the mass outside $[a,b]$ is ignored. It does not depend on $N$; it depends
on $L$.
\end{itemize}
The second one is more insidious than it looks. Below $a$ the cosine series represents, instead of the true payoff
$K-Fe^y$, its \emph{even reflection} about $a$, namely $K-Fe^{2a-y}$. The error is therefore
\begin{equation}\label{eq:cos-riflessione}
\varepsilon_{\text{tail}}\approx\int_{-\infty}^aF\big(e^{2a-y}-e^{y}\big)f(y)\dd y .
\end{equation}
The weight $e^{2a-y}$ \emph{grows} like $e^{-y}$ moving to the left. With $\rho<0$ the left tail of $X$ is heavy, so
even a tiny probability outside $[a,b]$ can produce a visible error. An interval based only on $c_2$ also ignores the
kurtosis.

\begin{esempio}[title={Example: the literature choice $L=12$}]
With $L=12$, $N=256$ (Fang--Oosterlee's choice for Heston), the maximum discrepancy between COS and Lewis on the test
grid is ${3.95\times10^{-5}}$ on the FO set and ${2.58\times10^{-6}}$ on the GS set (SPX-like): above the $10^{-6}$
threshold.

Diagnosis in three moves.
\begin{enumerate}
\item Increasing $N$ from 256 to 4096 does not change the error: it is not the series.
\item Increasing $L$ by 2 lowers the error by a factor of about 8: it is the interval.
\item On the FO set, $T=1$, ATM put, the term \eqref{eq:cos-riflessione} computed separately is
${3.948\times10^{-5}}$; the measured error is $-{3.947\times10^{-5}}$. The probability outside the interval is only
${7.7\times10^{-6}}$, but the weight $e^{-y}$ amplifies it.
\end{enumerate}
This is an instance of the principle ``suspect the math first, then the code, the test last''.
\end{esempio}

\subsection{Self-verifying COS}
A first fix was to set $L=18$. On the four test sets this was enough: maximum error ${1.01\times10^{-7}}$ (FO) and
${2.48\times10^{-9}}$ (GS). An independent check showed, however, that a fixed pair $(N,L)$ is not robust.

\begin{itemize}
\item On 300 random parameter sets plausible for the SPX, with 6 expiries each (1800 cells), COS with $(256,18)$
exceeds $10^{-6}$ in 413 cells (22.9\%), with a maximum of ${9.95\times10^{-2}}$.
\item The code issued a warning based on an estimate of the series error. But 5 cells were silently wrong (up to
${2.61\times10^{-6}}$), because the warning did not model the range error.
\end{itemize}

For one of these silent cells, Table~\ref{tab:cos-NL} shows the discrepancy from an independent reference price as
$(N,L)$ varies.

\begin{table}[htbp]
\centering
\small
\caption{Discrepancy of COS from an independent reference price, on a ``silent'' cell, as $(N,L)$ varies.}
\label{tab:cos-NL}
\begin{tabular}{@{}lccccc@{}}
\toprule
$(N,L)$ & $(256,18)$ & $(1024,18)$ & $(4096,18)$ & $(4096,24)$ & $(4096,30)$\\
\midrule
error & ${2.61\times10^{-6}}$ & ${2.51\times10^{-6}}$ & ${2.51\times10^{-6}}$ & ${2.09\times10^{-8}}$ & ${9.28\times10^{-10}}$\\
\bottomrule
\end{tabular}
\end{table}

Again: the error is flat in $N$ and collapses with $L$. In calibration (Chapter~\ref{cap:calibrazione}) the
parameters explore exactly these regions, so a COS that checks itself was needed.

\begin{idea}[title={Key idea: the self-verifying algorithm}]
\begin{enumerate}
\item Start from $(N,L)=(256,18)$ and compute the puts for all strikes of the expiry.
\item Recompute with $(2N,\,L+6)$ and measure the maximum discrepancy across strikes.
\item If the discrepancy is $\le10^{-9}F$ (that is, $10^{-7}$ at $F=100$, ten times below the threshold), return the
\emph{finer} price. Otherwise, double again.
\item If $N$ reaches $2^{14}$ without agreement, the code issues a ``not converged'' warning: it is never silently
wrong.
\end{enumerate}
Doubling $N$ squares the series error; adding 6 to $L$ reduces the range error by a factor of 20--500 (measured).
Hence the discrepancy between the two levels is a good estimate of the error of the coarser level, and the finer price
is better than the tolerance.
\end{idea}

\begin{attenzione}
Self-verification is an \emph{a posteriori} estimate, not a rigorous bound: it assumes that each refinement reduces
the error substantially. It also makes the price a piecewise function of the parameters (at the $10^{-9}F$ level),
which disturbs finite-difference Jacobians. This is why in calibration $(N,L)$ is frozen at a fixed value for the
whole fit.
\end{attenzione}

\begin{progetto}[title={In the qpricer project: self-verifying COS}]
\code{src/qpricer/engines/fourier.py}, function \code{cos\_price}: if $N$ and $L$ are not specified it uses the
self-verifying algorithm; if they are specified it performs a single evaluation (for tests and studies).
\begin{itemize}
\item Random scan (1800 cells): 0 cells above $10^{-6}$, 0 warnings, maximum error ${3.05\times10^{-9}}$.
\item Final $N$: 512 in 1248 cells, 1024 in 274, 2048 in 189, 4096 in 69, 8192 in 17, 16{,}384 in 3. Median cost 768
CF evaluations.
\item Time for 25 strikes: 0.30\,ms (GS, $T=1$), 0.62\,ms (FO, $T=1$), 7.97\,ms (worst cell, $T=3$), against 0.12\,ms
for fixed $(256,18)$ and 4.6--9.5\,ms for Lewis.
\end{itemize}
\end{progetto}

\section{Fourier pricing, method 2: the Lewis formula}
A second method, independent of the interval $[a,b]$ and of any series, serves as a check:
\begin{equation}\label{eq:lewis}
C=\DF\Big[F-\frac{\sqrt{FK}}{\pi}\int_0^\infty\mathrm{Re}\big[e^{iux}\varphi(u-\tfrac i2)\big]\frac{\mathrm{d}u}{u^2+\frac14}\Big],
\qquad x=\ln(F/K).
\end{equation}
The put is obtained by replacing the first $F$ with $K$.

\subsection{Step-by-step derivation}
\textbf{Step 1.} $(F_T-K)^+=F_T-\min(F_T,K)$ and $\E F_T=F$. It is therefore enough to compute $\E\min(F_T,K)$.

\textbf{Step 2.} With $F_T=Fe^X$ we write $\min(F_T,K)=K\,g(x+X)$, where $g(y)=\min(e^y,1)$.

\textbf{Step 3.} The Fourier transform of $g$ exists in the strip $0<\mathrm{Im}\,\zeta<1$. Splitting the integral
into $y<0$ (where $g=e^y$) and $y>0$ (where $g=1$):
\[
\hat g(\zeta)=\int_{\R}e^{i\zeta y}g(y)\dd y=\frac{1}{1+i\zeta}-\frac{1}{i\zeta}=\frac{1}{\zeta^2-i\zeta}.
\]
\textbf{Step 4.} The transform is inverted along the line $\zeta=u+\tfrac i2$, in the middle of the strip:
\[
\E\min(F_T,K)=\frac{K}{2\pi}\int_{\R}\hat g(\zeta)\,e^{-i\zeta x}\,\varphi(-\zeta)\dd u .
\]
On this line $\zeta^2-i\zeta=u^2+\tfrac14$ is real and positive, and $Ke^{x/2}=\sqrt{FK}$.

\textbf{Step 5.} By the symmetry $\varphi(-u-\tfrac i2)=\overline{\varphi(u-\tfrac i2)}$ the integral over $\R$ folds
onto $[0,\infty)$ as a real part: this gives \eqref{eq:lewis}.

\begin{perche}[title={\texorpdfstring{Why the contour $\mathrm{Im}\,u=-\tfrac12$ in particular}{Why the contour Im u = -1/2 in particular}}]
The contour only requires $\E[F_T^{1/2}]<\infty$, which always holds (by Jensen's inequality
$\E[F_T^{1/2}]\le F^{1/2}$). Other contours (for example those of Carr--Madan) require higher moments, which in Heston
can explode (Section~\ref{sec:heston-code}). Moreover, the integrand is regular at $u=0$: we have
$4\varphi(-\tfrac i2)\le4$, without the $1/u$ singularity of the Gil-Pelaez formulas.

Check with Black: if $X$ is Gaussian with total variance $s^2$, $\varphi(u-\tfrac i2)=e^{-(u^2+1/4)s^2/2}$.
\end{perche}

\subsection{The quadrature}
The integral is computed with adaptive quadrature on $[0,\infty)$, vectorized over all strikes of an expiry. The price
error is at most $\DF\sqrt{FK}\cdot\mathrm{err}/\pi$, where $\mathrm{err}$ is the quadrature error.

Keeping it below $10^{-7}$ up to $K=150$ requires $\mathrm{err}\le{2.6\times10^{-9}}$. With SciPy's default
tolerances (${1.49\times10^{-8}}$) the price error reaches about ${2\times10^{-6}}$, which is not enough: the code uses
tolerances of $10^{-11}$ and raises an error if the quadrature does not report convergence.

\begin{progetto}
\code{lewis\_price} in \code{src/qpricer/engines/fourier.py}. In the checks of this chapter the quadrature never
failed: on 1800 random cells, on all the extreme cases tried, and for $T$ down to $10^{-4}$. Against an independent
reference price the agreement is $\le{1.62\times10^{-9}}$ on the test grid.
\end{progetto}

\section{\texorpdfstring{The limit $\xi\to0$ and the correction term $\rho\xi J$}{The limit xi -> 0 and the correction term rho xi J}}
\subsection{Zeroth order: Black at the mean variance}
If the vol-of-vol is zero, the variance is deterministic ($v_t=\bar v(t)$) and $F_T$ is lognormal with total variance
$\bar\Sigma$. The price is therefore Black's with $\sigma_{\text{eff}}^2=\bar\Sigma/T$.

It seems natural to check the code like this: ``at $\xi=10^{-4}$ Heston must agree with Black to within $10^{-5}$''.
But with $\rho\ne0$ this is not true, and it is instructive to understand why.

\subsection{\texorpdfstring{First order in $\xi$}{First order in xi}}
\textbf{Step 1.} Expand the solution of the Riccati equation: $D=D_0+\xi D_1+O(\xi^2)$. At zeroth order
$D_0=-\tfrac12A\,\tau_\kappa(\tau)$; at first order $D_1'=-\kappa D_1+i\rho uD_0$.

\textbf{Step 2.} Integrating, and using $C=\kappa\theta\int D$:
\[
\ln\varphi(u)=-\tfrac12A\,\bar\Sigma+\xi\rho\,(iu)\big(-\tfrac12A\big)J+O(\xi^2).
\]
\textbf{Step 3.} The Black CF with total variance $\Sigma$ is $e^{-A\Sigma/2}$. A factor $-\tfrac12A$ therefore acts as
a derivative with respect to $\Sigma$, and a factor $iu$ as a derivative with respect to $x=\ln F$ (integrating by
parts in the price). We obtain
\begin{equation}\label{eq:xi-zero}
V_H=V_B(\sigma_{\text{eff}})+\rho\,\xi\,J\,\partial_x\partial_\Sigma V_B+O(\xi^2),
\end{equation}
with
\begin{align}
J&=\int_0^T\bar v(t)\,\frac{1-e^{-\kappa(T-t)}}{\kappa}\dd t
=\frac{\theta}{\kappa}(T-\tau_\kappa)+\frac{v_0-\theta}{\kappa}(\tau_\kappa-TE),\label{eq:J}\\
\partial_x\partial_\Sigma V_B&=-\frac{\DF\,F\,n(d_1)\,d_2}{2\bar\Sigma}.\nonumber
\end{align}
Here $n$ is the standard normal density and $d_{1,2}$ are Black's at $\sigma_{\text{eff}}$.

\begin{idea}
With correlation $\rho\ne0$ the distance from Black is \emph{linear} in $\xi$, not quadratic. The gate ``at
$\xi=10^{-4}$ Heston agrees with Black to within $10^{-5}$'' is therefore achievable only for $\rho=0$. For $\rho\ne0$
the correct reference is Black \emph{plus} the term \eqref{eq:xi-zero}.
\end{idea}

\begin{esempio}[title={Example: how large the correction is}]
$\kappa=1.5$, $v_0=\theta=0.04$, $\rho=-0.7$, $\xi=10^{-4}$, $T=1$, $K=120$, $F=100$, $\DF=1$. The hand estimate of
the term $\rho\xi J\partial_x\partial_\Sigma V_B$ is $-{3.3\times10^{-4}}$: thirty times the $10^{-5}$ tolerance.

An independent numerical check (Lewis against Black, same $\kappa,\theta,T$) confirms the linearity: the maximum
discrepancy is ${3.27\times10^{-3}}$, ${3.27\times10^{-4}}$, ${3.25\times10^{-5}}$ for $\xi=10^{-3},10^{-4},10^{-5}$. At
$\rho=0$ and $\xi=10^{-4}$ it is instead ${2.6\times10^{-8}}$.
\end{esempio}

The adopted solution is a two-part gate, with the tolerance unchanged at $10^{-5}$:
\begin{itemize}
\item at $\rho=0$ against pure Black at $\sigma_{\text{eff}}$;
\item at $\rho=-0.7$ against Black plus the term \eqref{eq:xi-zero}.
\end{itemize}
Reducing $\xi$ to $10^{-7}$ to make the linear term vanish was rejected: it would test the rounding in the
$\kappa\theta/\xi^2$ terms, not the model.

\begin{progetto}[title={\texorpdfstring{In the qpricer project: the limit $\xi\to0$ measured}{In the qpricer project: the limit xi -> 0 measured}}]
Grid $K\in[60,150]$, $T\in[0.1,3]$, $\xi=10^{-4}$, $\kappa=1.5$, both engines (COS and Lewis), threshold $10^{-5}$:
\begin{center}
\small
\begin{tabular}{@{}lrrr@{}}
\toprule
 & $\rho=0$ & $\rho=-0.7$ & $\rho=-0.7$\\
$(\theta,v_0)$ & vs Black & vs Black${}+{}$\eqref{eq:xi-zero} & vs pure Black\\
\midrule
$(0.04,\,0.04)$ & ${4.36\times10^{-8}}$ & ${5.58\times10^{-8}}$ & ${6.05\times10^{-4}}$\\
$(0.04,\,0.09)$ & ${4.18\times10^{-8}}$ & ${5.33\times10^{-8}}$ & ${6.58\times10^{-4}}$\\
$(0.09,\,0.01)$ & ${3.02\times10^{-8}}$ & ${3.66\times10^{-8}}$ & ${6.20\times10^{-4}}$\\
\bottomrule
\end{tabular}
\end{center}
The cases with $v_0\ne\theta$ also check the time averaging in $\sigma_{\text{eff}}$. An independent check
differentiated its own price numerically at $\xi=0$: the first-order coefficient agrees with $\rho J\partial_x
\partial_\Sigma V_B$ to within ${1.6\times10^{-10}}$. The gap from pure Black therefore lies in the model, not in the
code.
\end{progetto}

As $T\to0$ we have $J\to v_0T^2/2$, and \eqref{eq:xi-zero} becomes a shift in implied volatility:
\begin{equation}\label{eq:heston-skew}
\delta\sigma\approx\frac{\rho\,\xi}{4\sqrt{v_0}}\ln\frac KF,
\qquad\text{that is, ATM skew}\approx\frac{\rho\,\xi}{4\sqrt{v_0}} .
\end{equation}
This is why $\rho<0$ produces the skew. It also reveals a limitation of the model: as $T\to0$ the skew tends to a
\emph{finite} value, whereas the very short-dated SPX skew keeps steepening.

\section{Tails, moment explosion, symmetry}
\label{sec:heston-code}
\paragraph{Moment explosion.} With large $\xi$ some moments $\E[(F_T/F)^\omega]$ become infinite after a finite time
$T^*(\omega)$. For $\rho\le0$, with $\chi=\rho\xi\omega-\kappa$ and $\Delta_\omega=\chi^2-\xi^2(\omega^2-\omega)$:
\begin{equation}\label{eq:esplosione}
T^*(\omega)=\begin{cases}
\infty & \Delta_\omega\ge0,\\[2pt]
\dfrac{2}{\sqrt{-\Delta_\omega}}\Big[\arctan\dfrac{\sqrt{-\Delta_\omega}}{\chi}+\pi\,\ind_{\chi<0}\Big] & \Delta_\omega<0.
\end{cases}
\end{equation}
No $\omega\in[0,1]$ ever explodes. For example, on the FO set $T^*(-2)=4.28$ years (hand estimate).

Two practical consequences.
\begin{itemize}
\item The slope of the smile wings is fixed by the exploding moments (Lee's moment formula).
\item A Fourier contour that requires an exploded moment diverges. Lewis's contour, which requires only
$\omega=\tfrac12$, is always safe.
\end{itemize}

\paragraph{Symmetry.} Changing measure with density $F_T/F$ (the share measure), Heston remains Heston with modified
parameters. By Girsanov the variance acquires the drift $\rho\xi v$, and $Y=-X$ has correlation $-\rho$:
\[
C(F,K;\kappa,\theta,\xi,\rho,v_0)=\frac KF\,P\Big(F,\frac{F^2}{K};\ \kappa-\rho\xi,\ \frac{\kappa\theta}{\kappa-\rho\xi},
\ \xi,\ -\rho,\ v_0\Big).
\]
It is a good test: it links a call and a put computed with different parameters.

\section{Numerical pitfalls}
\begin{enumerate}
\item \textbf{Original form of the CF}: branch cut of the logarithm (wrong prices at long $T$) and overflow of
$e^{dT}$. Use the little trap form.
\item \textbf{Catastrophic cancellation in $\beta-d$} for small $\xi$, and the $2\kappa\theta/\xi^2$ amplification of the
logarithm's error: use \eqref{eq:heston-cf-code} with \code{log1p}.
\item \textbf{Calls computed directly with COS}: errors amplified by $e^b$. Compute puts.
\item \textbf{Interval too narrow}: the error stops decreasing with $N$ but changes with $L$. Use self-verifying COS.
\item \textbf{Strikes outside $[a,b]$}: without truncating $z$ one gets wrong prices, even negative ones. The resulting
zero prices for deep OTM strikes must be treated as missing in the volatility inversion.
\item \textbf{The $c_2$ typo} in Fang--Oosterlee: use \eqref{eq:c2}.
\item \textbf{Lewis tolerances}: the quadrature defaults are not enough.
\item \textbf{$\rho\to\pm1$}: the CF decays only like $e^{-c\sqrt u}$ and COS convergence becomes algebraic; we
restrict $|\rho|\le0.99$.
\item \textbf{Feller violated}: not a problem for Fourier; it is one for Monte Carlo.
\end{enumerate}

\section{Checks and published values}
The project uses four parameter sets for its checks:
\begin{itemize}
\item \textbf{FO}: $\kappa=1.5768$, $\theta=0.0398$, $\xi=0.5751$, $\rho=-0.5711$, $v_0=0.0175$; a Eurostoxx-50
calibration used in the literature; violates Feller ($\nu_F=0.38$);
\item \textbf{GS}: $\kappa=1.3253$, $\theta=0.0354$, $\xi=0.3877$, $\rho=-0.7165$, $v_0=\theta$; SPX-like
(Gatheral 2006); violates Feller ($\nu_F=0.62$);
\item \textbf{LW}: $\kappa=4$, $\theta=0.25$, $\xi=1$, $\rho=-0.5$, $v_0=0.04$ ($\nu_F=2$);
\item \textbf{AT2}: $\kappa=1.5$, $\theta=v_0=0.04$, $\xi=0.3$, $\rho=-0.9$ ($\nu_F=1.33$).
\end{itemize}

\begin{progetto}[title={In the qpricer project: checks of the Fourier part}]
\textbf{COS against Lewis} on $K\in[60,150]$, $T\in[0.1,3]$, $\DF\in\{1,0.95\}$, calls and puts, threshold
$10^{-6}$:
\begin{center}
\small
\begin{tabular}{@{}lrrrr@{}}
\toprule
 & FO & GS & LW & AT2\\
\midrule
self-verifying COS & ${1.83\times10^{-10}}$ & ${1.40\times10^{-12}}$ & ${1.14\times10^{-13}}$ & ${5.12\times10^{-12}}$\\
fixed $(256,18)$ & ${1.01\times10^{-7}}$ & ${2.48\times10^{-9}}$ & ${1.10\times10^{-13}}$ & ${9.82\times10^{-9}}$\\
fixed $(256,12)$ & ${3.95\times10^{-5}}$ & ${2.58\times10^{-6}}$ & ${8.35\times10^{-9}}$ & ${1.36\times10^{-7}}$\\
\bottomrule
\end{tabular}
\end{center}

\textbf{Published values} (project discrepancy):
\begin{itemize}
\item Fang--Oosterlee, FO set, ATM, $T=1$: 5.785155450, reproduced to $-{1.6\times10^{-8}}$; at $T=10$:
22.318945791, to $+{1.5\times10^{-10}}$.
\item Lewis reference prices (via QuantLib): to ${3\times10^{-14}}$.
\item Broadie--Kaya (4 published decimals): to ${4.2\times10^{-5}}$.
\item Andersen's set A-I, $K=100$, value 13.08467014 (single source): Lewis to $-{3.0\times10^{-9}}$.
\end{itemize}
The values of Albrecher et al.\ (2007) are \emph{not} reproduced to the stated precision for 3 values out of 8 (for
example $+0.0073$ at $T=2$ with stated precision 0.005). No single rate explains them all; the most likely diagnosis is
that the parameters are printed to 4 digits, and rounding them moves the prices by more than the stated precision.
For this reason those values are not used as binding tests; the trap itself is checked by the negative control.
\end{progetto}

\begin{riassunto}
\begin{itemize}
\item The SPX skew corresponds to a heavy left tail, that is, to the leverage effect. Heston: CIR variance correlated
with the price.
\item $v_0,\theta$ levels; $\kappa$ speed; $\xi$ curvature; $\rho$ slope. Feller matters for MC, not for Fourier.
\item Closed-form CF from Feynman--Kac, the affine ansatz and Riccati. We use the little trap form, written without
subtractions and with \code{log1p}.
\item COS: density as a cosine series, coefficients from the CF, exponential convergence; puts are priced. The range
error is the real risk: the code uses a self-verifying $(N,L)$.
\item Lewis: a single integral on the contour $\mathrm{Im}\,u=-\frac12$, always defined; it is the independent check.
\item $\xi\to0$: Black at the mean variance only for $\rho=0$; otherwise there is a term linear in $\rho\xi$.
\end{itemize}
\end{riassunto}

\needspace{8\baselineskip}
\begin{domande}
\begin{enumerate}
\item What effect does each Heston parameter have on the smile?
\item Why does the Heston CF have a closed form? \emph{Hint:} affine model, exponential-affine ansatz, Riccati.
\item What is the ``little Heston trap''? Why do two algebraically equal formulas give different prices?
\item Why does COS price puts rather than calls?
\item The COS error does not decrease when $N$ is increased. What do you conclude, and what do you change?
\item Why does self-verification compare $(N,L)$ with $(2N,L+6)$ and return the finer price?
\item The Feller condition is violated: is the COS price wrong? \emph{Hint:} no; it matters for simulation.
\item Why, with $\rho\ne0$, does the limit $\xi\to0$ not agree with Black to within $10^{-5}$?
\end{enumerate}
\end{domande}
